\section{Memory Models and Atomic Registers}

\subsection{Memory Reordering and Hierarchy}

\begin{remark}
    In a concurrent environment, programs often fail even if they appear logically sound in a sequential context. This is primarily due to \textbf{Memory Reordering}. Compilers and hardware are permitted to change the order of memory operations as long as the semantics of a \textbf{sequentially executed program} remain unaffected.
\end{remark}

Memory reordering offers significant performance benefits through optimizations like dead code elimination, register hoisting, and locality optimization, but it introduces errors when multi-threaded consistency is assumed.

\begin{codeexample}
int x;
// Thread A: Waiting for a signal
void wait() { 
    x = 1; 
    while (x == 1); 
}
// Thread B: Providing the signal
void arrive() { 
    x = 2; 
}
\end{codeexample}

Consider how the \textbf{wait()} function might be compiled:
\textbf{Naive Assembly (Standard):} The value of \textbf{x} is re-read from memory in every iteration.

\begin{codeexample}
movl 0x1, x  
test:
mov x, %eax    // Load x from memory
cmp 0x1, %eax // Compare x to 1
je test        // Jump back if equal
\end{codeexample}

\textbf{Optimized Assembly (The "Failure" Case):} The compiler assumes \textbf{x} cannot change within the loop (since Thread A doesn't modify it there) and hoists the value into a register once or replaces the check with an infinite jump.

\begin{codeexample}
movl 0x1, x
test:
jmp test       // Infinite loop: Thread A never sees Thread B's update!
\end{codeexample}

\begin{definition} 
    \textbf{Memory Hierarchy}:
    \begin{itemize} 
        \item \textbf{Registers:} Fast, low latency (0.5ns), but very low capacity.
        \item \textbf{L1/L2 Caches:} Intermediate speed and capacity (1ns - 7ns). 
        \item \textbf{System Memory:} Slow, high latency (100ns), but high capacity.
    \end{itemize}
\end{definition}

\begin{important}
    \textbf{Hardware View on Reordering \& Visibility:} Modern CPUs heavily rely on \textbf{Pipelines}, allowing instructions to be reorganized internally (out-of-order) to prevent stalling. Furthermore, because L1 caches and registers are \textbf{private to each core}, a \texttt{store} operation to memory on Core 1 might be buffered locally and only become visible to Core 2 at a later, unpredictable time.
\end{important}

\begin{remark}
    \textbf{Architecture Models Table:} The lecture compares memory models across architectures (x86, ARM, Alpha). A crucial takeaway is that almost all common architectures (except conceptually strict ones) generally allow the reordering of \textbf{Stores after Loads} to maximize performance, unless explicitly prevented by memory barriers.
\end{remark}

\subsection{The Java Memory Model (JMM)}

\begin{definition} 
    A \textbf{Memory Model} is a contract between the programmer and the system. It provides the minimal guarantees for memory operation effects while allowing hardware/compiler optimizations.
\end{definition}

\begin{definition} 
    \textbf{Program Order (PO):} A total order of actions within a \textbf{single} thread as defined by the execution trace. Crucially, PO does \textbf{not} provide a global ordering guarantee across threads.
\end{definition}

\begin{definition} 
    \textbf{Synchronization Actions (SA):} Actions like reading/writing \textbf{volatile} variables or acquiring/releasing locks.
\end{definition}

\begin{definition} 
    \textbf{Synchronization Order (SO):} A total order of all SA. All threads see these actions in the same global order.
\end{definition}

\begin{definition} 
    \textbf{Happens-Before (HB):} The transitive closure of PO and \textbf{Synchronizes-With} (SW). If action $A$ happens-before $B$, $B$ is guaranteed to see the results of $A$.
\end{definition}

\begin{remark}
    \textbf{HB-Consistency \& Data Races:} The JMM explicitly allows "Data Races" (when two threads access the same variable without HB-ordering and at least one is a write). In an \textbf{HB-consistent execution}, a read operation is allowed to see either the temporally last write ordered before it via the HB-relation, \textit{or} any other write that is not strictly ordered by HB. This is why unsynchronized reads can observe stale or surprising values.
\end{remark}

\subsection{Registers and Concurrent Executions}

\begin{important}
    \textbf{Important Distinction:} In the context of this theory, a \textbf{Register} does \textit{not} refer to a hardware CPU register (like \texttt{\%eax}), but rather to a fundamental, abstract memory object that provides \texttt{read()} and \texttt{write(v)} operations.
\end{important}

\begin{definition}
    A Register can act with different consistency properties: \textbf{Safe}, \textbf{Regular}, or \textbf{Atomic}. Every operation $J$ on a register has an \textbf{effect time} $\tau(J)$ that lies somewhere between the invocation (start) and response (end) of the operation.
\end{definition}

\begin{theorem}
    \textbf{Safe Register:} A register where a read not overlapping with any write returns the correct written value. If a read overlaps with a write, it can return \textit{any} value from the variable's domain.
\end{theorem}

\begin{theorem}
    \textbf{Regular Register:} Like a safe register, but if a read overlaps with a write, it returns either the old value or the new value being written (no random garbage data is returned).
\end{theorem}

\begin{theorem}
    \textbf{Atomic Register:} The strongest guarantee. Any read or write operation $J$ appears to happen instantaneously at its specific effect time $\tau(J)$ between its invocation and completion. 
    \begin{itemize}
        \item Two operations on the same register always have different effect times (no true simultaneity).
        \item A \texttt{read()} always returns the value of the temporally closest preceding \texttt{write()} based on these effect times.
    \end{itemize}
\end{theorem}

\subsection{Solving Mutual Exclusion}

To implement a Lock (Mutex) correctly using only basic memory reads/writes, three conditions must be met:
\begin{enumerate}
    \item \textbf{Mutual Exclusion:} No two processes can be in the critical section (CS) simultaneously.
    \item \textbf{Freedom from Deadlock:} If processes try to enter the CS, at least one must eventually succeed.
    \item \textbf{Freedom from Starvation:} Every process that attempts to enter must eventually succeed.
\end{enumerate}

\begin{remark}
    \textbf{Combinatorics of Interleavings:} For 2 threads executing $k$ statements each, there are exactly $\binom{2k}{k} = O\left(\frac{4^k}{\sqrt{2k}}\right)$ possible interleavings.
    Because this state space explodes so rapidly, manual "proof by exhaustion" through simple trace analysis is impossible. This strongly motivates the use of systematic \textbf{State Space Diagrams}. Furthermore, the lecture uses this explosion to justify \textbf{reduced state space diagrams}, which simplify the analysis by grouping or omitting states that aren't strictly structurally relevant to the synchronization protocol itself.
\end{remark}

\begin{definition}
    A \textbf{State Space Diagram} is a formal tool used to track every possible configuration of a concurrent program.
    \textbf{States} are typically represented as a tuple, e.g., $[p, q, \text{wantp}, \text{wantq}]$, where $p$ and $q$ are the current program counters (line numbers) of the threads.
\end{definition}

\begin{important}
    \textbf{Failure Identification} in State Space Diagrams:
    \begin{itemize}
        \item \textbf{No Mutex:} Reachable state where both processes are in the CS.
        \item \textbf{Deadlock:} A state with no outgoing transitions where threads are stuck in the protocol.
        \item \textbf{Starvation:} An execution path where one thread never reaches the CS despite trying.
    \end{itemize}
\end{important}

\subsection{Historical "Failed" Attempts}

Before Peterson's solution, several intuitive but flawed approaches were developed for two processes $P$ and $Q$.

\begin{enumerate}
    \item \textbf{Check then Set (The "After You" Problem)}
    \begin{codeexample}
// Process P (Process Q is symmetric)
while(wantq);
wantp = true;
// Critical Section
wantp = false;
    \end{codeexample}
    \textbf{Failure:} \textbf{Mutual Exclusion fails}. If both check \textbf{want} simultaneously, see \textbf{false}, both set to \textbf{true} and enter the CS.

    \item \textbf{Set then Check (The Deadlock Problem)}
    \begin{codeexample}
// Process P
wantp = true;
while(wantq);
// Critical Section
wantp = false;
    \end{codeexample}
    \textbf{Failure:} \textbf{Deadlock}. If both set flags simultaneously, they wait forever for the other.

    \item \textbf{Strict Alternation (The Starvation Problem)}
    \begin{codeexample}
// Process P
while(turn != 1);
// Critical Section
turn = 2;
    \end{codeexample}
    \textbf{Failure:} \textbf{Starvation}. If Process Q never wants to enter, Process P can enter once, set \textbf{turn = 2}, and then wait forever because Q won't set it back to 1.
\end{enumerate}

\subsection{Dekker's Algorithm}
The first correct solution to the mutual exclusion problem for two processes. It combines the flag system with a \textbf{turn} variable to break ties. If both want to enter, the \textbf{turn} variable decides who must temporarily retract their interest.

\begin{codeexample}
// Process P
wantp = true;
while (wantq) {
    if (turn == 2) {
        wantp = false;
        while (turn == 2); // wait for turn
        wantp = true;
    }
}
// Critical Section
turn = 2;
wantp = false;
\end{codeexample}

\begin{theorem}
    Dekker's algorithm guarantees Mutual Exclusion, Freedom from Deadlock, and Freedom from Starvation.
\end{theorem}

\subsection{Peterson's Algorithm (The Modern Standard)}
Peterson's algorithm is a more concise version of Dekker's. Instead of retracting interest, a process simply "gives way" by setting a \textbf{victim} variable to itself.

\begin{codeexample}
// Process P
wantp = true;
victim = 1; // P is process 1
while (wantq && victim == 1);
// Critical Section
wantp = false;
\end{codeexample}

\begin{important}
    \textbf{Wait Condition:} A process waits only if the other is interested \textbf{AND} it is currently its own turn to be the "victim".
\end{important}

\begin{theorem}
    Peterson's Algorithm provides Mutual Exclusion, Deadlock Freedom, and Starvation Freedom, and works elegantly using only atomic registers.
\end{theorem}

\begin{proof}
    \textbf{Mutual Exclusion:} We assume atomic registers where $W_i(x)$ is a write to $x$ and $R_i(x)$ is a read of $x$ by process $i$.
    Assume concurrent critical sections $CS_P$ and $CS_Q$. Assume without loss of generality (w.l.o.g.):
    $$W_Q(\text{victim} = Q) \rightarrow W_P(\text{victim} = P)$$

    From the algorithm's code for process $P$:
    $$W_P(\text{flag}[P]=\text{true}) \rightarrow W_P(\text{victim}=P) \rightarrow R_P(\text{flag}[Q]) \rightarrow R_P(\text{victim}) \rightarrow CS_P$$

    To enter $CS_P$, process $P$ must observe the while condition as false:
    \begin{itemize}
        \item \textbf{Observation 1:} From (1), $W_P(\text{victim}=P)$ is the latest write. Thus, $P$ must read \texttt{victim} $= P$.
        \item \textbf{Observation 2:} Since $P$ reads \texttt{victim} $= P$, it must read \texttt{flag}[Q] $= \text{false}$ to enter $CS_P$.
    \end{itemize}
    From the code for $Q$ and the transitivity of the "$\rightarrow$" relation:
    $$W_Q(\text{flag}[Q]=\text{true}) \rightarrow W_Q(\text{victim}=Q) \rightarrow W_P(\text{victim}=P) \rightarrow R_P(\text{flag}[Q])$$ 
    This implies $P$ must read \texttt{flag}[Q] $= \text{true}$. 

    The requirement to read false (Observation 2) contradicts the forced read of true (from Equation 3). Thus, concurrent execution is impossible.
\end{proof}

\begin{important}
    \textbf{Java Implementation Warning:} A naive implementation of Peterson's in Java might use a \texttt{volatile boolean[]} array for the \texttt{want} flags. However, in Java, \texttt{volatile} on an array \textit{only protects the array reference itself}, not the individual array contents! To implement this correctly and guarantee visibility across threads, one must use an \texttt{AtomicIntegerArray}.
\end{important}